% Standalone audit: compile directly with pdfLaTeX, twice.
% MAIN VERDICT: MINOR REVISION
% CORE VERDICT: MINOR REVISION
% EXTENDED VERDICT: VERIFIED
% P02 v1.0 and all supplied project files remain unchanged.
\documentclass[11pt]{article}
\usepackage[T1]{fontenc}
\usepackage[utf8]{inputenc}
\usepackage{lmodern,amsmath,amssymb,amsthm}
\usepackage[margin=0.9in]{geometry}
\usepackage{array,booktabs,longtable,enumitem,needspace}
\usepackage[hidelinks]{hyperref}
\newcolumntype{P}[1]{>{\raggedright\arraybackslash}p{#1}}
\setlength{\emergencystretch}{3em}
\setlength{\parindent}{0pt}
\setlength{\parskip}{5pt}
\setlist{nosep,leftmargin=*}
\newcommand{\RR}{\mathbb R}
\newcommand{\ZZ}{\mathbb Z}
\newcommand{\one}{\mathbf 1}
\newcommand{\Dom}{(-1,\infty)^n}
\newcommand{\Orth}{\RR_{\geq0}^{n}}
\newcommand{\lin}[1]{\textbf{L#1}}
\newcommand{\pass}{\textbf{PASS}}
\newcommand{\sha}[2]{\par{\small\ttfamily #1\allowbreak #2}\par}
\title{A02: Independent mathematical audit of P02\\
Synchronous subspace and equilibria}
\author{Independent reconstruction and adversarial review}
\date{27 September 2026}
\begin{document}
\maketitle

\section{Verdict and exact version}

\textbf{Main verdict: MINOR REVISION.}
The canonical invariance, scalar reduction, complete root classification,
positive-equilibrium theorem, domain filtering, algebraic loci, and
multiplicity statements are mathematically correct. One local wording
correction is required in the boundary-continuity paragraph:
\lin{545--552} introduces \(s_\pm\), then supplies limits
\emph{respectively} in the order \((s_-,s_+)\).
Explicitly naming that ordered pair removes a branch-label ambiguity.
This is issue E01 below. No displayed root formula, count, assumption,
theorem conclusion, or intended limit needs alteration.

\textbf{CORE VERDICT: MINOR REVISION.}\quad
\textbf{EXTENDED VERDICT: VERIFIED.}
The extended verdict concerns only the limited claims of
\lin{45--48,440--453}: invariance for homogeneous finite real coefficients
and the positive-root alternatives for \(A>0\). It does not export the
canonical negative-root table or global nonnegative-flow conclusions
to arbitrary signed coefficients.

The exact audited source is
\path{proofs/P02_synchronous_subspace_and_equilibria.tex},
\textbf{P02 v1.0, dated 2026-09-27, DRAFT}, 26,167 bytes and 598 lines,
from the supplied \path{cyclic_network_project_P02_draft.zip}.
Its full-file SHA-256, including comments and whitespace, is
\sha{a63f0757eb0de6462b7ce18c179f6f5c21}{079a8c3e8d0f9c3ed53d6ce7b256f3}
The ZIP's SHA-256 is
\sha{68e5b8a2277ee89d37816a440e9c487b36}{b2ec873aea336fba02ad83bd6eaa7e}
All P02 line references in A02 refer to that exact 598-line file,
including blank and comment lines. This is a new standalone audit;
P02, the manifest, registry, compendium, and other supplied project
files have not been revised or promoted.

The supplied ZIP was selected deliberately. Two pre-existing local
project folders contain a 1,754-byte P02 placeholder with a different
hash; neither was used as the proof under review. The ZIP's current
manifest identifies P02 v1.0 with the matching digest above. Its current
registry and compendium agree that P01 is audited, P02 is DRAFT, and
P03--P18 remain TODO. Historical all-TODO prose in planning documents
is explicitly superseded by current dated records.

The audit starts from the canonical equation and standard local ODE
uniqueness. It does not rely on compilation, internal check totals,
archived derivations, or later modules. No literature search, numerical
sampling as proof, simulation, P02 rewrite, or P03 work was performed.
The exact finite checks below only corroborate the analytic proofs.

\section{Current inputs and exact assumptions}

All nine requested current files were located inside the same archive
and used. Their full hashes are included in the accompanying inventory.
\begin{longtable}{@{}P{0.50\textwidth}P{0.44\textwidth}@{}}
\toprule
Current file & Use in the audit\\
\midrule
\endhead
\path{project_manifest.md} & Current version/hash, P01 closure, status precedence.\\
\path{claim_registry.md} & C025, C026, C028--C033, C146 and P02-D1--D4.\\
\path{frozen_theorem_package.md} & S2 scope, next P03 attempt, separate later modules.\\
\path{00_canonical_model.tex} & Exact field, domain, parameters, indexing.\\
\path{01_notation_and_assumptions.tex} & A1--A3, synchrony, state sets, abbreviations.\\
\path{proofs/P01_wellposedness_and_absorbing_box.tex} & Audited uniqueness and canonical global nonnegative continuation.\\
\path{audits/A01_wellposedness_and_absorbing_box.tex} & Audit and sole P01 correction E01.\\
\path{proofs/P02_synchronous_subspace_and_equilibria.tex} & Exact draft under independent review.\\
\path{proofs_compendium.tex} & Current banner and actual P02 input; no proof premise.\\
\bottomrule
\end{longtable}

The core assumptions are exactly
\begin{equation}\label{A02:assumptions}
n\in\ZZ,\ n\ge5,\quad 1\le q\le n-1,\qquad
a,k,c,g,h>0,\qquad U=\Dom.
\end{equation}
Indices lie in \(\ZZ/n\ZZ\), and all terms remain when offsets coincide.
Use only as notation
\[
A=k+h,\qquad \mu=a+c-g,\qquad b=a-A.
\]
Thus \(A>0\) and \(b>-A\) follow from the assumptions; neither \(\mu\)
nor \(b\) has an assumed sign. The local \(b\) is not an archived
comparator coefficient. No coprimality condition is imposed.
The only vector field audited is
\begin{equation}\label{A02:field}
F_i(x)=ax_i-kx_i^2+\frac{cx_{i-1}}{1+x_i}
-\frac{gx_i}{1+x_{i+q}}-hx_ix_{i-q}.
\end{equation}
The nonnegative model state space is \(\Orth\).
Negative synchronous equilibria in \(U\) belong to the open extension,
not the nonnegative model.

\section{Independent invariance and scalar reconstruction}

For \(x=s\one\), \(s>-1\), each coordinate appearing in every component
equals \(s\). Direct substitution gives, for every \(i\),
\[
F_i(s\one)=as-ks^2+\frac{cs}{1+s}
-\frac{gs}{1+s}-hs^2
=as-As^2+\frac{(c-g)s}{1+s}.
\]
Thus \(F(s\one)=f(s)\one\). This uses no \(n,q\), gcd, cycle
connectedness, or independence of neighboring coordinates. In particular
\(q=1\), \(q=n-1\), and \(2q=0\pmod n\) cause no change.
The same substitution works for \(n=1\), any residue \(q\), and finite
real homogeneous coefficients.

Both fields are smooth on their respective open domains. A scalar
solution lifts to a vector solution; uniqueness identifies it with the
vector solution from the same synchronous initial point. If the scalar
maximal interval ended at a finite time while the vector solution
continued in \(U\), continuity would give a synchronous limit
\(s_*\one\in U\). Local scalar existence at \(s_*>-1\) would extend
the scalar solution, a contradiction. Conversely a longer scalar
interval would extend the vector solution. The maximal intervals
therefore coincide. The argument also applies backward locally.
This proves invariance throughout the maximal interval in \(U\);
it does not assert global existence for negative initial values.

The inherited scalar domain is exactly \(s>-1\). On it,
\begin{align}
f(s)&=s\left[a-As+\frac{c-g}{1+s}\right]\notag\\
&=\frac{s}{1+s}\left[a+(a-A)s-As^2+c-g\right]\notag\\
&=\frac{s}{1+s}\left[\mu+bs-As^2\right].
\label{A02:scalar}
\end{align}
At \(g=c\), cancellation along synchrony does not enlarge the declared
domain: the original rational terms remain undefined at \(s=-1\).
If one considered the rational expression outside \(U\), the same
tangency argument would work on each interval avoiding \(-1\);
P02's equilibrium classification concerns the specified \(U\).
For \(s(0)\ge0\), P01 supplies a global nonnegative vector solution,
and invariance gives a global nonnegative scalar solution. No
absorbing-radius, spectral, stability, or later-module result is needed.

\section{Origin, reversible reduction, explicit roots, and substitution}

All five original terms vanish at \(s=0\), so the origin is retained
before division. For \(s\ne0\), \(s>-1\),
\[
f(s)=0
\Longleftrightarrow \mu+bs-As^2=0
\Longleftrightarrow Q(s):=As^2-bs-\mu=0.
\]
Division by \(s\) and multiplication by \(1+s>0\) are reversible
on precisely this set. Hence the synchronous equilibrium set in \(U\) is
\begin{equation}\label{A02:set}
\{0\one\}\ \cup\ \{s\one:Q(s)=0,\ s>-1,\ s\ne0\}.
\end{equation}
Zero roots of \(Q\) do not add another origin.

Completing the square independently yields
\begin{align}
Q(s)&=A\left(s-\frac b{2A}\right)^2-\frac{\Delta}{4A},\notag\\
\Delta&=b^2+4A\mu=(a-A)^2+4A(a+c-g)
=(a+A)^2+4A(c-g).
\label{A02:disc}
\end{align}
The draft's \(D\) is this \(\Delta\). Since \(A>0\), there are
no real roots for \(\Delta<0\), one double root for \(\Delta=0\),
and two distinct simple roots for \(\Delta>0\).
For \(\Delta\ge0\), keep the fixed labels
\begin{equation}\label{A02:roots}
s_-=\frac{b-\sqrt{\Delta}}{2A}\le
s_+=\frac{b+\sqrt{\Delta}}{2A}.
\end{equation}
Direct calculation, and independently coefficient comparison, give
\begin{equation}\label{A02:vieta}
s_-+s_+=\frac bA,\qquad
s_-s_+=\frac{b^2-\Delta}{4A^2}=-\frac{\mu}{A}.
\end{equation}
For \(r=(b+\varepsilon\sqrt{\Delta})/(2A)\), \(\varepsilon=\pm1\),
\[
\mu+br-Ar^2=\mu+\frac{b^2-\Delta}{4A}=0.
\]
Substitution into the original scalar expression, not just \(Q\), gives
\begin{equation}\label{A02:residual}
ar-kr^2+\frac{cr}{1+r}-\frac{gr}{1+r}-hr^2
=\frac{r}{1+r}(\mu+br-Ar^2)=0\qquad(r\ne-1).
\end{equation}
Only \(r>-1\) is retained as an equilibrium in \(U\).
For \(r<-1\) this is an identity for the rational expression outside
the declared domain, not a DM-domain equilibrium claim.
The special formulas at \(\mu=0,\Delta=0,g=c,b=0\), and the exact
near-zero branch are all covered by this direct substitution.
A truncated branch expansion is not an exact root.

The singular-root test is
\begin{equation}\label{A02:minusone}
Q(-1)=A+b-\mu=a-(a+c-g)=g-c.
\end{equation}
Thus \(-1\) is a polynomial root if and only if \(g=c\).
It is never admissible. Neither multiplication by \(1+s\) nor
cancellation of equal rational terms admits that point.

\section{Complete canonical classification}

\subsection{Domain filtering}
Whenever \(\Delta\ge0\),
\[
s_+\ge\frac b{2A}=\frac a{2A}-\frac12>-\frac12>-1.
\]
For the smaller root, since \(a+A>0\),
\begin{align}
s_->-1
&\Longleftrightarrow \sqrt{\Delta}<b+2A=a+A\notag\\
&\Longleftrightarrow \Delta<(a+A)^2
\Longleftrightarrow g>c.
\label{A02:domain}
\end{align}
Squaring is legitimate because both sides are nonnegative and the
right side is strictly positive. Equality gives \(s_-=-1\) exactly
when \(g=c\); reversing the inequalities gives \(s_-<-1\) exactly
when \(g<c\). These statements are conditional on real roots existing.
At \(g=c\),
\[
Q(s)=(s+1)(As-a),
\]
so the distinct roots are exactly \(-1,a/A\).

\subsection{Every sign of \texorpdfstring{\(\mu\) and \(b\)}{mu and b}}
If \(\mu>0\), then \(\Delta>b^2\), so \(s_-<0<s_+\), both simple,
for each of \(b<0,b=0,b>0\). The negative root is below \(-1\),
equal to \(-1\), or in \((-1,0)\) according as \(g<c,g=c,g>c\).
In the last cell one also has \(g<a+c\).

If \(\mu=0\), \(Q=s(As-b)\). For \(b>0\), the roots are
\(s_-=0,s_+=b/A>0\). For \(b<0\), they are
\(-1<s_-=b/A<0=s_+\), using \(b>-A\). Both are simple for \(b\ne0\).
For \(b=0\), zero is the single distinct polynomial root, of
multiplicity two. In all these subcases \(g=a+c>c\).

If \(\mu<0\), then \(g>a+c>c\), so every real root is admissible.
If \(b=0\), \(\Delta=4A\mu<0\), and there are no real roots.
For \(b\ne0,\Delta\ge0\), the positive Vieta product and the sign of
the sum show that both nonzero roots have the sign of \(b\).
They are distinct and simple for \(\Delta>0\), coincident and double
for \(\Delta=0\). For \(b<0\), the double root is in \((-1/2,0)\).
For either nonzero sign of \(b\), \(\Delta<0\) gives no real roots.
Thus all three signs of \(b\) are included in the no-real-root cell.

An asterisk in this complete table marks one distinct polynomial root
of multiplicity two; every other root displayed is simple.
\begin{longtable}{@{}c P{0.25\textwidth}P{0.36\textwidth}P{0.23\textwidth}@{}}
\toprule
Cell & Conditions & Polynomial roots and signs & Roots in \(s>-1\)\\
\midrule
\endhead
I & \(\mu>0,\ g<c\) & \(s_-<-1<0<s_+\) & \(s_+\) only\\
II & \(\mu>0,\ g=c\) & \(s_-=-1,\ s_+=a/A>0\) & \(s_+\) only\\
III & \(\mu>0,\ g>c\) & \(-1<s_-<0<s_+\) & Both\\
IV & \(\mu=0,\ b>0\) & \(s_-=0,\ s_+=b/A>0\) & Both\\
V & \(\mu=0,\ b=0\) & \(s_-=s_+=0^*\) & The one zero\\
VI & \(\mu=0,\ b<0\) & \(-1<s_-=b/A<0=s_+\) & Both\\
VII & \(\mu<0,\ \Delta<0\) & None; any sign of \(b\) & None\\
VIII & \(\mu<0,\ b>0,\ \Delta>0\) & \(0<s_-<s_+\) & Both\\
IX & \(\mu<0,\ b>0,\ \Delta=0\) & \(s_-=s_+=(b/(2A))^*>0\) & One distinct root\\
X & \(\mu<0,\ b<0,\ \Delta>0\) & \(-1<s_-<s_+<0\) & Both\\
XI & \(\mu<0,\ b<0,\ \Delta=0\) & \(-1/2<s_-=s_+=(b/(2A))^*<0\) & One distinct root\\
\bottomrule
\end{longtable}

Let \(R_Q\) count distinct real polynomial roots including zero,
\(D_Q\) those in the domain, and \(E_U,E_{\ge0},E_+\) all synchronous
equilibria in \(U\), nonnegative equilibria, and positive equilibria.
The last column uses P02's definitions: \(r,d\) count only nonzero
polynomial and nonzero admissible roots.
\begin{center}
\begin{tabular}{@{}c c c c c c c@{}}
\toprule
Cell & \(R_Q\)&\(D_Q\)&\(E_U\)&\(E_{\ge0}\)&\(E_+\)&P02 \((r,d,N_0,N_+)\)\\
\midrule
I&2&1&2&2&1&\((2,1,2,1)\)\\
II&2&1&2&2&1&\((2,1,2,1)\)\\
III&2&2&3&2&1&\((2,2,2,1)\)\\
IV&2&2&2&2&1&\((1,1,2,1)\)\\
V&1&1&1&1&0&\((0,0,1,0)\)\\
VI&2&2&2&1&0&\((1,1,1,0)\)\\
VII&0&0&1&1&0&\((0,0,1,0)\)\\
VIII&2&2&3&3&2&\((2,2,3,2)\)\\
IX&1&1&2&2&1&\((1,1,2,1)\)\\
X&2&2&3&1&0&\((2,2,1,0)\)\\
XI&1&1&2&1&0&\((1,1,1,0)\)\\
\bottomrule
\end{tabular}
\end{center}
Every P02 table entry passes. The zero root of \(Q\) is nonnegative;
negative roots never are. Every positive polynomial root is admissible.
The origin is added once, whether or not it is a root of \(Q\).

\section{Positive-equilibrium theorem: both directions and equalities}

Two distinct positive roots imply by Vieta that \(\mu<0\), \(b>0\),
and \(\Delta>0\). Conversely those three conditions give two distinct
real roots with positive product and positive sum, hence both positive.
All positive roots satisfy \(s>-1\). Define
\begin{equation}\label{A02:gf}
g_f=a+c+\frac{(a-A)^2}{4A}
=c+\frac{(a+A)^2}{4A},\qquad
s_f=\frac{a-A}{2A}.
\end{equation}
Since \(\Delta=4A(g_f-g)\), necessity and sufficiency give exactly
\begin{equation}\label{A02:two}
E_+=2\quad\Longleftrightarrow\quad
a>A,\qquad a+c<g<g_f.
\end{equation}
None of these strict inequalities can be weakened to include an endpoint.

The full positive classification is
\begin{center}
\begin{tabular}{@{}lll@{}}
\toprule
Relation & Parameter range & Positive roots\\
\midrule
\(a>A\)& \(0<g\le a+c\)& One simple root \(s_+\)\\
\(a>A\)& \(a+c<g<g_f\)& Two simple roots \(s_-,s_+\)\\
\(a>A\)& \(g=g_f\)& One double root \(s_f>0\)\\
\(a>A\)& \(g>g_f\)& None\\
\(a\le A\)& \(0<g<a+c\)& One simple root \(s_+\)\\
\(a\le A\)& \(g\ge a+c\)& None\\
\bottomrule
\end{tabular}
\end{center}
At \(a>A,g=a+c\), \(Q\) has roots \(0,b/A\), the latter positive
and simple. At \(a=A,g=a+c\), \(Q\) has double zero, not a positive
root. At \(a<A,g=a+c\), the other root is negative.
At \(a<A,g=g_f\), the double root is negative; at
\(a=A,g=g_f=a+c\), it is zero. Thus a repeated root is not
unconditionally positive.

An independent positive-half-line argument uses
\(H(s)=\mu+bs-As^2\). If \(b\le0\), \(H'(s)=b-2As<0\)
for \(s>0\), so precisely \(\mu>0\) yields one positive zero.
If \(b>0\), the maximum at \(b/(2A)>0\) equals \(\Delta/(4A)\).
For \(\mu<0\), a positive, zero, or negative maximum gives two,
one double, or no positive zeros. For \(\mu=0\), factor
\(H=s(b-As)\); for \(\mu>0\), the negative Vieta product leaves
one positive root. This proves every implication in P02's four
alternatives, including no-positive and exact-one-distinct-root clauses.

\section{Negative admissible roots and the singular boundary}

The complete negative-root conclusions follow independently from
Vieta and \eqref{A02:domain}:
\begin{enumerate}
\item For \(\mu>0\), exactly one polynomial root is negative.
It is below \(-1\) for \(g<c\), equals \(-1\) for \(g=c\), and
lies in \((-1,0)\) for \(c<g<a+c\).
\item For \(\mu=0\), exactly one negative root exists if and only if
\(a<A\). It is \(a/A-1\in(-1,0)\). Otherwise there is no negative root.
\item For \(\mu<0,a<A\), there are two distinct negative roots for
\(\Delta>0\), one double negative root for \(\Delta=0\), and none
for \(\Delta<0\). All real ones lie in \((-1,0)\).
\item For \(\mu<0,a\ge A\), there are no negative roots:
at \(a=A\) there are no real roots; at \(a>A\) the real roots,
if any, are positive.
\end{enumerate}
For \(a<A,\mu<0\), these conditions can also be written
\(a+c<g<g_f\) for two negative roots, \(g=g_f\) for one double,
and \(g>g_f\) for none. Both roots below \(-1\), or a larger root
at \(-1\), are impossible canonically. Indeed \(s_+>-1/2\);
also \(g\le c\) forces \(\mu\ge a>0\), so if one root is at or
below \(-1\), the other is positive.

Thus one or both polynomial roots can be domain-admissible negative
roots exactly in the cases listed, but neither represents a
nonnegative state. The smaller root alone crosses the singular boundary
at \(g=c\), where \(\Delta=(a+A)^2>0\). Excluding negative states
from the physical model does not justify dropping them from a
classification on \(U\).

\section{Algebraic loci, expansions, and exact zero multiplicities}

Solving \(Q=0\) for \(g\) is a reversible linear rearrangement:
\begin{equation}\label{A02:graph}
g=a+c+bs-As^2=g_f-A(s-s_f)^2.
\end{equation}
For nonzero canonical equilibria, retain \(s>-1,s\ne0,g>0\).
The algebraic continuation includes the origin exactly at \(g=a+c\).
The nonzero branch itself does not literally contain the origin;
P02's qualification by algebraic continuation is correct.

A repeated root obeys \(Q=Q'=0\). Since \(Q'=2As-b\), it is
\(s_f=b/(2A)\), and substitution gives \(\mu=-b^2/(4A)\), or
\(g=g_f\). Conversely those values make \(Q=A(s-s_f)^2\).
Canonically \(s_f>-1/2\), and \(s_f>0\) if and only if \(a>A\).
The two loci meet exactly when
\(g_f-(a+c)=b^2/(4A)=0\), that is \(a=A,g=a+c\).
These conclusions are algebraic, not bifurcation theorems.

At a nonzero root in the domain, the factor \(-s/(1+s)\) in
\(f=-sQ/(1+s)\) is smooth and nonzero, so \(f\) and \(Q\)
have the same root multiplicity there.
At zero, the numerator is \(\mu s+bs^2-As^3\) and the denominator
does not vanish:
\begin{center}
\begin{tabular}{@{}lll@{}}
\toprule
Parameters & Multiplicity of zero in \(f\)& Multiplicity of zero in \(Q\)\\
\midrule
\(\mu\ne0\)& One & Not a root\\
\(\mu=0,b\ne0\)& Two & One\\
\(\mu=b=0\)& Three & Two\\
\bottomrule
\end{tabular}
\end{center}
Direct substitution at \(a=A,g=a+c\) gives
\[
f(s)=As-As^2-\frac{As}{1+s}=-\frac{As^3}{1+s},
\qquad Q(s)=As^2.
\]
Because \(A>0\), these orders are exactly three and two, and zero
is the only synchronous equilibrium.
P02 \lin{343--345} explicitly attaches the triple zero to \(f\)
and the double root to \(Q\). Its degeneracy terminology is
unambiguous; no correction is required.

For fixed \(b\ne0\), the square-root expansion at \(\mu=0\) is
\[
\sqrt{b^2+4A\mu}
=|b|\left(1+\frac{2A\mu}{b^2}
-\frac{2A^2\mu^2}{b^4}+O(\mu^3)\right).
\]
Therefore the root tending to zero is
\[
s_0(\mu)=\frac{b-\operatorname{sgn}(b)\sqrt{b^2+4A\mu}}{2A}
=-\frac{\mu}{b}+\frac{A\mu^2}{b^3}+O(\mu^3).
\]
It is \(s_-\) when \(b>0\), and \(s_+\) when \(b<0\).
For sufficiently small nonzero \(\mu\), this root is nonzero and
domain-admissible, and \(g=a+c-\mu>0\). At \(\mu=0\),
\[
f(s)=\frac{s^2(b-As)}{1+s}
=bs^2-(A+b)s^3+O(s^4)=bs^2-as^3+O(s^4).
\]
Both draft expansions and the quadratic coefficient \(b\) are correct.
They assert no full-system normal form.

When \(b=0\), \(Q=As^2-\mu\). For \(\mu>0\),
\(s_\pm=\pm\sqrt{\mu/A}\), and the negative root is admitted exactly
when \(0<\mu<A\), equivalently \(c<g<a+c\). It is \(-1\) at
\(\mu=A\), below \(-1\) at \(\mu>A\). At \(\mu=0\) both labels
give zero; at \(\mu<0\) there are no real roots. This also passes.

\section{Boundary continuity and issue E01}

The pair in \eqref{A02:roots} is jointly continuous on the relative
parameter set \(A>0,\Delta\ge0\). Always in the explicit increasing
order \((s_-,s_+)\), the boundary limits are
\begin{align*}
g\to a+c,\ b>0:\quad &(s_-,s_+)\to(0,b/A),\\
g\to a+c,\ b<0:\quad &(s_-,s_+)\to(b/A,0),\\
g\to g_f,\ \Delta\downarrow0:\quad &(s_-,s_+)\to(s_f,s_f),\\
g\to c:\quad &(s_-,s_+)\to(-1,a/A).
\end{align*}
The last boundary retains only \(a/A\) as a nonzero domain equilibrium.
For \(\Delta<0\), there are no real branches to continue.

At a point with \(a=A,\mu>0\), continuity gives
\((-\sqrt{\mu/A},\sqrt{\mu/A})\). At \(a=A,\mu<0\), the
discriminant is strictly negative, so a whole neighborhood has no
real roots. At the simultaneous degeneracy \(b=\mu=0\), every
approach through \(\Delta\ge0\) to a canonical point with \(A>0\)
has \(b,\sqrt{\Delta}\to0\), so both labels tend to zero.
For fixed \(b=0\), this is the one-sided approach \(\mu\downarrow0\).
Along \(\mu=0\) with \(b\) varying,
\[
s_-=\min\{b/A,0\},\qquad s_+=\max\{b/A,0\}.
\]
Both are continuous; the identity of the label equal to zero changes
only where the roots actually coalesce. The exact formulas and table
never relabel branches.

\textbf{E01 (minor, required wording clarification).}
P02 \lin{545} introduces the ordered formulas \(s_\pm\).
At \lin{548--549}, the limits given \emph{respectively} as
\(0,b/A\) for \(b>0\), and \(b/A,0\) for \(b<0\), use the order
\((s_-,s_+)\). The antecedent \(s_\pm\) can instead be read in
its conventional expansion \(s_+,s_-\).
The word \emph{ordered} suggests the intended increasing order but
does not explicitly name it. This is a localized ambiguity, not
a mathematical failure of any formula or intended limit.

For an exact diagnostic, fix \(k=h=1/2,A=1,a=2,c=1\) and
let \(g\to3\). Then \(b=1,\mu\to0\), and
\[
(s_-,s_+)\to(0,1),\qquad (s_+,s_-)\to(1,0).
\]
The second pair cannot have the limits assigned by a literal
plus-then-minus reading of the draft sentence.
Earlier P02 \lin{187,206--208,372--373} already fixes the correct labels.

\textbf{Minimal proposed correction, not applied to P02:}
replace the implicit pairing at \lin{548--552} with
\begin{quote}
At \(\mu\to0\) with fixed \(b>0\),
\((s_-,s_+)\to(0,b/A)\); with fixed \(b<0\),
\((s_-,s_+)\to(b/A,0)\). With \(b=0\) and \(\mu\downarrow0\),
both labels tend to zero. At \(g\to c\),
\((s_-,s_+)\to(-1,a/A)\); the first root is excluded from the
chosen domain at the boundary.
\end{quote}
Alternatively, explicitly declaring the increasing pair at the
paragraph's start resolves the same issue. No root should be renamed.
Because branch continuity was an explicit audit requirement,
this clarification is required for an unqualified pass of the exact text.

\section{Domain reversibility of every transformation}

\begin{longtable}{@{}P{0.26\textwidth}P{0.66\textwidth}@{}}
\toprule
Operation & Reversibility and limitation\\
\midrule
\endhead
Factoring \(s\) & An identity for every \(s>-1\), including zero.
The original field is checked at zero first.\\
Dividing by \(s\) & Only for \(s\ne0\).
At \(\mu=0\), zero roots of \(Q\) are excluded from nonzero counts.\\
Multiplying by \(1+s\) & Reversible for \(s>-1\).
The polynomial root \(-1\) at \(g=c\) is discarded.\\
Changing to \(Q\) & Multiplication by \(-1\) changes no zeros;
the nonzero and domain filters remain.\\
Completing the square & Exact with \(A>0\); the three discriminant
signs exhaust every possible real-root count.\\
Taking square roots & Only for \(\Delta\ge0\), retaining both signs
and counting a repeated value once.\\
Squaring the domain test & Uses \(\Delta\ge0\) and \(a+A>0\);
not extended to unrestricted signed parameters.\\
Solving for \(g\) & Linear rearrangement without division by \(s\).
The continuation graph must still be filtered by \(s>-1,s\ne0,g>0\).\\
Taylor expansions & Local formulas with remainders at fixed \(b\ne0\)
or near \(s=0\); truncations are not claimed to be exact roots.\\
\bottomrule
\end{longtable}
No polynomial-only claim is incorrectly retained as an equilibrium
of the original rational vector field. The excluded \(-1\), the
already counted zero, and roots outside \(U\) are consistently separated.

\section{Optional weaker-assumption material}

For finite real homogeneous coefficients, any integer \(n\ge1\),
and any residue \(q\), substitution and smoothness/local uniqueness
prove synchronous invariance on the domain. If only \(A=k+h>0\)
is imposed, with real \(a,c,g\), the polynomial, discriminant, roots,
and Vieta identities remain valid. On \(s>0\), the factor
\(-s/(1+s)\) is nonzero. The sign/vertex proof of the four positive
alternatives uses only \(A>0\); the positive-root extension is valid.

The negative-domain table does not follow under those assumptions.
The domain comparison needed \(a+A>0\), and the larger-root and
repeated-root lower bounds used \(a>0\). P02 states these restrictions
explicitly and makes no classification claim for \(A=0\) or \(A<0\).

The following exact signed tests use \(A=1\):
\begin{longtable}{@{}P{0.30\textwidth}P{0.62\textwidth}@{}}
\toprule
\((a,c,g)\) & Exact consequence\\
\midrule
\endhead
\((-1,-2,-4)\) & \(\mu=1,Q=s^2+2s-1\).
Roots \(-1\pm\sqrt2\) give one positive root.\\
\((2,-10,-63/8)\) & \(\mu=-1/8,\Delta=1/2\);
two distinct positive roots.\\
\((2,-10,-31/4)\) & \(\mu=-1/4,Q=(s-1/2)^2\);
one double positive root.\\
\((2,-10,-15/2)\) & \(\mu=-1/2,\Delta=-1\); no real roots.\\
\((-1,-2,-3)\) & \(\mu=0,Q=s(s+2)\); no positive root,
and \(-2\) is outside the domain.\\
\((1,-2,-1)\) & \(\mu=b=0,Q=s^2\); origin triple for \(f\), double for \(Q\).\\
\((-3,1,1)\) & \(g=c,Q=(s+3)(s+1)\); both roots excluded.
Here the larger root is \(-1\).\\
\((-3,1,2)\) & \(g>c,Q=(s+2)^2\); repeated root outside the domain.\\
\((-1,1,1)\) & \(a+A=0,Q=(s+1)^2\); repeated root is the excluded \(-1\).\\
\bottomrule
\end{longtable}
Taking \(k=-2,h=3\) in the first example additionally tests that the
extended scalar conclusions use the sum, not separate signs.
For each retained root the original five-term residual is exactly
\eqref{A02:residual}. The last three examples refute only an
\emph{unclaimed} extension of the canonical negative-root table.
They do not challenge P02's limited extension.
\textbf{The separate extended verdict is VERIFIED.}

\section{Adversarial counterexample-search record}

\subsection{Independently chosen canonical representatives}
In this table \(k=h=1/2,A=1,c=1\); every tuple is strictly canonical.
Roots are exact, not rounded. A star denotes a double polynomial
root. The origin is an equilibrium in every row and is counted once.
\begin{longtable}{@{}P{0.28\textwidth}c c P{0.33\textwidth}@{}}
\toprule
Required regime & \(a\)&\(g\)& Roots or obstruction\\
\midrule
\endhead
\(a<A,\mu>0\)&\(1/2\)&\(5/4\)&
\((-1-\sqrt5)/4,(-1+\sqrt5)/4\), both in domain\\
\(a=A,\mu>0\)&\(1\)&\(7/4\)& \(-1/2,1/2\)\\
\(a>A,\mu>0\)&\(2\)&\(9/4\)& \(-1/2,3/2\)\\
\(a<A,\mu=0\)&\(1/2\)&\(3/2\)& \(-1/2,0\)\\
\(a=A,\mu=0\)&\(1\)&\(2\)& \(0^*\)\\
\(a>A,\mu=0\)&\(2\)&\(3\)& \(0,1\)\\
\(a<A,\mu<0,\Delta>0\)&\(1/2\)&\(99/64\)& \(-3/8,-1/8\)\\
\(a<A,\mu<0,\Delta=0\)&\(1/2\)&\(25/16\)& \((-1/4)^*\)\\
\(a<A,\mu<0,\Delta<0\)&\(1/2\)&\(2\)& \(\Delta=-7/4\), no real roots\\
\(a=A,\mu<0\)&\(1\)&\(9/4\)& \(\Delta=-1\), no real roots\\
\(a>A,\mu<0,\Delta>0\)&\(2\)&\(51/16\)& \(1/4,3/4\)\\
\(a>A,\mu<0,\Delta=0\)&\(2\)&\(13/4\)& \((1/2)^*\)\\
\(a>A,\mu<0,\Delta<0\)&\(2\)&\(7/2\)& \(\Delta=-1\), no real roots\\
\(g=c\)&\(2\)&\(1\)& \(-1,2\); exclude \(-1\)\\
\(g<c\)&\(1/2\)&\(1/2\)&
\((-1-\sqrt{17})/4<-1,\ (-1+\sqrt{17})/4>0\)\\
\bottomrule
\end{longtable}
This includes the entire requested minimum adversarial list, all
signs of \(b\) for \(\mu>0\) and \(\mu=0\), and all possible
discriminant subcases for \(\mu<0\).
Exact comparisons \(2<\sqrt5<3\) and \(4<\sqrt{17}<5\)
locate the irrational roots without decimal evaluation.

For \(a=2,g=51/16\),
\[
Q=(s-1/4)(s-3/4),\qquad a+c=3<51/16.
\]
The original scalar residuals, grouping the two quadratic terms and
the two rational terms only for display, are
\[
f(1/4)=\frac12-\frac1{16}-\frac7{16}=0,\qquad
f(3/4)=\frac32-\frac9{16}-\frac{15}{16}=0.
\]
This breaks the archived claim that \(a+c\le g\) precludes positive
synchrony. The equality witness \(a=2,g=3,s=1\) breaks it at
equality as well. P02 correctly rejects that archived statement.

For \(a=1/2,g=99/64\),
\[
Q=(s+3/8)(s+1/8),
\]
and the original residuals are
\[
f(-3/8)=-\frac3{16}-\frac9{64}+\frac{21}{64}=0,\qquad
f(-1/8)=-\frac1{16}-\frac1{64}+\frac5{64}=0.
\]
Two admissible nonzero roots therefore need not be positive.
The negative double-root example refutes an unqualified positive
interpretation of the repeated-root locus. The \(g=c\) example tests
the canceled-denominator trap. The case \(a=A,\mu=0\) tests both
double-counting and the distinction between triple and double zeros.
No counterexample to a correctly interpreted P02 theorem was found;
the proofs above, rather than absence of a counterexample, establish
its mathematical conclusions.

\subsection{Independent reproduction of the finite-check record}

All eleven tuples at P02 \lin{505--515} give the printed
\(\mu,\Delta,(r,d,N_0,N_+)\). The additional choices at
\lin{521--524} are valid: for each \(a=1/2,1,2\), choosing
\(g=1/2,1,1+a/2\) realizes cells I, II, III respectively.
Choosing \(g=g_f+1\) gives cell VII for all three signs of \(b\).

The set at \lin{529--542} was independently reconstructed with
\[
V=\{1/4,1/2,1,2,3,4\},\qquad
\eta=\min\{a,A,c,1\}/64,
\]
using \(V^4\) and the three neighboring values of each
\(g=c,a+c,g_f\) for every \((a,A,c)\in V^3\).
Its union has exactly \textbf{2,868} distinct positive tuples.
Counting each distinct real root once per tuple and evaluating the
eleven printed examples separately gives exactly \textbf{4,438}
defined original scalar residuals, all zero, and \textbf{217}
occurrences of the undefined root \(-1\), all excluded.
A double root counts as one value. Nine of the eleven examples
already belong to the grid; these totals count additional evaluations,
not 2,879 distinct tuples.

The included script evaluates the original five-term rational
expression using exact arithmetic \(u+v\sqrt{\Delta}\).
Signs are decided by rational comparisons, with no floating-point
approximation. It checks discriminant identities, Vieta, domain
comparisons, and both directions of the positive alternatives.
Fifteen additional independently chosen representatives above pass.
These finite checks corroborate the exhaustive analytic audit and
verify the reproducible check totals. Other historical operation counts
reported only in the manifest are not treated as proof or claimed to
have been recovered from unprovided original scripts.

\section{Dependency verification and scope discipline}

P01 v1.1's full-file SHA-256 is
\sha{0d6950821aa83f3d09e8bfe84dab2b0b5}{020aa6e7881cf6928e984bff1890f1e}
\Needspace{3\baselineskip}
Its body hash, from the first P01 section line through EOF, is
\sha{a640b2bf9a160feed273dc667c6b44a692}{3a1d68d6dd81acd35640cc0e4d5808}
The full-file hash matches the current manifest; the body hash matches
both the manifest and P01 header.
The sole A01 correction is present at current P01 line 436: the
archived radius is described as the larger real root and as nonnegative.
A01 correctly retains its historical verdict for its earlier audited
hash; the current closure record identifies the corrected version.

P02 uses only P01's local smoothness and uniqueness
(\path{prop:P01-local}, current P01 \lin{80--93}) and canonical
global nonnegative continuation (\path{thm:P01-global}).
Its equilibrium algebra does not need the absorbing-box proof.
P01's local statement already allows finite real coefficients;
standard uniqueness for the time-reversed local equation also justifies
the two-sided local invariance assertion.

No premise from P03 or later mathematics enters any proof.
Including later placeholders in the compendium is not a mathematical
dependency. The frozen package's S2 role is respected: P02 classifies
synchronous equilibria only. It asserts no classification of all
nonsynchronous equilibria, stability, bistability, transcritical
bifurcation, saddle-node/fold bifurcation, Hopf bifurcation,
or Fourier spectral conclusion.

The multiplicity adjectives are read in the scalar sense explicitly
fixed at P02 \lin{202--205,255--264}. A simple scalar zero does
not establish full-vector Jacobian nonsingularity or a simple
eigenvalue. The branch loci and scalar quadratic coefficient are
algebraic inputs for possible later work, not completed bifurcation
results. This contextual qualification is sufficient; no additional
P02 correction is required.

\section{Line-by-line findings and issue list}

The ledger covers every substantive P02 paragraph. Blank lines,
environment delimiters, and labels within each range carry no
additional claim. E01 is the sole required correction.
\begin{longtable}{@{}P{0.14\textwidth}P{0.16\textwidth}P{0.62\textwidth}@{}}
\toprule
P02 lines & Finding & Independent reason or action\\
\midrule
\endhead
1--25 & \pass & Version, assumptions, P01 status and scope agree with current metadata.\\
30--48 & \pass & Invariance and inherited domain; no gcd or distinct-offset condition;
finite-real-coefficient extension valid.\\
51--75 & \pass & Component substitution, uniqueness and continuation prove
invariance on the full maximal interval.\\
79--95 & \pass & \(b>-A\), scalar identities, canonical global nonnegative
inference, and unchanged domain at \(g=c\).\\
99--123 & \pass & Origin, exact polynomial, discriminant identities and both root formulas.\\
125--150 & \pass & Reversible operations, completed square, singular root test,
and no additional zero equilibrium.\\
155--174 & \pass & Definitions and distinct-value counting; origin included once.\\
184--200 & \pass & All eleven cells and counts; no missing \(b=0\) case.\\
202--209 & \pass & Scalar multiplicities and correct zero-branch labels;
negative roots excluded only from nonnegative states.\\
211--234 & \pass & Vieta and domain test with both sign conditions for squaring.\\
236--253 & \pass & Exhaustive \(\mu\) trichotomy; real roots for \(\mu<0\)
admitted since \(g>a+c>c\).\\
255--265 & \pass & Multiplicity transfer and origin orders, algebraic continuation,
no stability inference.\\
269--294 & \pass & Four mutually exclusive positive alternatives and every converse;
maximum two and exact-one distinction.\\
296--309 & \pass & Independent monotonicity/vertex argument is valid.\\
313--347 & \pass & \(g_f,s_f\), loci and their intersection;
repeated-root sign/domain and unambiguous triple/double distinction.\\
349--366 & \pass & Reversible graph in \(g\), correct filters and degeneracy.\\
368--385 & \pass & Near-zero root label and both Taylor coefficients correct;
cubic coefficient is \(-a\).\\
387--398 & \pass & Square-root branches and singular threshold; bifurcation
work explicitly reserved.\\
402--438 & \pass & Original-parameter iff classification and all equalities;
topology independence.\\
440--453 & \pass & Positive extension requires only \(A>0\); negative-domain
classification is correctly not exported.\\
458--491 & \pass & Original scalar residuals and specializations,
discriminant expansion and Vieta identities.\\
493--518 & \pass & All eleven rational examples, discriminants and counts checked.\\
519--527 & \pass & Equality/sign representative choices valid;
finite checks do not promote status.\\
529--542 & \pass & Independently reproduced 2,868 tuples, 4,438 defined
residuals and 217 excluded singular-root occurrences.\\
544--547 & \pass & Continuity on \(A>0,\Delta\ge0\) and coalescence correct.\\
545,548--552 & \textbf{E01: minor} & Explicitly name \((s_-,s_+)\)
when assigning the listed limits. Intended limits and formulas are correct.\\
550--555 & \pass & Both-zero and \(b\to0\) limits correct;
domain counts may change at a boundary.\\
558--573 & \pass & Completeness and adversarial distinctions correctly retain
singular, zero, negative and repeated cases.\\
575--598 & \pass & Checklist is reference material, not task authority;
all summarized mathematical points checked. Internal checks do not replace audit.\\
\bottomrule
\end{longtable}

\clearpage
\section{Downstream consequences and revision closure}

The underlying canonical results are established by this independent
reconstruction with fixed labels \eqref{A02:roots}. The exact P02 v1.0
text nevertheless receives \textbf{MINOR REVISION}, rather than an
unqualified VERIFIED label, because E01 leaves a branch-order ambiguity
in one sentence. No project status metadata was changed.

Closure requires the local E01 clarification in a separately authorized
P02 revision, recording its new hash, and checking that the branch
labels and other mathematical content are preserved. No change to
the canonical model, P01, displayed root formulas, classification
tables, proofs, or optional extension is required. A02 remains
tied to the exact v1.0 hash recorded above.

\textbf{P03 is cleared to begin after E01 is corrected and closure is
recorded.} No other P02 mathematical blocker remains.
P03 has not been begun here. Later full-system linearization must
still be derived independently with coincident-offset terms retained.
Two positive roots do not establish bistability; the algebraic loci
do not establish bifurcation type; scalar simplicity does not certify
full-system nonsingularity. No P03--P18 claim is promoted by A02.

The audit-only package contains this standalone source, a PDF reading
copy, exact input hashes, a reproducibility script and its result,
and a short provenance/closure note. It contains no replacement
P02, canonical source, registry, or compendium.

\end{document}
